\begin{pastexam}[label={px:ncu-106-dsa-07}]{中央資工所 106 學年度 第 7 題}
（問答）

Below is the famous Dijkstra Shortest Path Algorithm (DSPA). The input of DSPA is a weighted
directed graph $G=(V, E)$ and a specific source node $s$, where $V$ is the node set and $E$ is the
edge set and every edge in $E$ has a positive weight. (Note that the weight of the edge $(u, v)$ is
stored in $ew[u, v]$). For every node $u$ in $V$-$\{s\}$, DSPA can derive the shortest path from
$s$ to $u$. However, DSPA does not consider node weights. If we now consider node weight and
include node weights in the total weight of every path, we need to extend the original DSAP.
Please extend DSAP to consider the node weights. Note that we assume the weight of a node $u$ is
stored in $nw[u]$ and the weight of the destination node is not included in the total weight of
the path. Note that DSPA uses indentation to represent the block structure. So please use
indentation properly when writing down the extended DSPA algorithm. You should write down the
complete algorithm including the input, the output and all steps. (10\%)

\begin{center}
\small
\begin{tabular}{|p{0.92\linewidth}|}
\hline
\textbf{Algorithm}: DSPA (Dijkstra Shortest Path Algorithm)\\ \hline
\textbf{Input}: $G=(V, E)$, $ew$, $s$ //$G=(V, E)$ is a graph with edge weights stored in $ew$, and $s$ is the source\\
\textbf{Output}: $SP$ //$SP$ is the set of shortest paths from $s$ to all other nodes\\ \hline
1:\ \ $dist[s]\leftarrow 0$; $dist[u]\leftarrow\infty$, for each $u\neq s$, $u\in V$\\
2:\ \ \textbf{insert} $u$ with key $dist[u]$ into priority queue $Q$, for each $u\in V$\\
3:\ \ \textbf{while} ($Q\neq\varnothing$)\\
4:\ \ \hspace*{1.5em}$u\leftarrow$Extract-Min($Q$)\\
5:\ \ \hspace*{1.5em}\textbf{for} each $v$ adjacent to $u$\\
6:\ \ \hspace*{3em}\textbf{if} $dist[v] > dist[u]+ew[u,v]$ \textbf{then}\\
7:\ \ \hspace*{4.5em}$dist[v]\leftarrow dist[u]+ew[u,v]$\\
8:\ \ \hspace*{4.5em}$pred[v]\leftarrow u$\quad //$v$'s predecessor in the shortest path is $u$\\
9:\ \ \textbf{calculate} the shortest path from $s$ to $u$ to add into set $SP$ according to $pred[u]$, for each $u\in V$, $u\neq s$\\
10: \textbf{return} $SP$\\ \hline
\end{tabular}
\end{center}
\end{pastexam}

\begin{pastsoln}
題目裡的 DSAP 是 DSPA 的筆誤。 題目沒有說節點權重的正負，這裡假設 $nw[u] \ge 0$（理由見最後）。

\textbf{路徑的權重}。 路徑 $P = \langle s = v_0, v_1, \ldots, v_k = t\rangle$ 的總權重是
\[
w(P) = \sum_{i=0}^{k-1} ew[v_i, v_{i+1}] + \sum_{i=0}^{k-1} nw[v_i],
\]
節點權重算到 $v_{k-1}$，終點 $t$ 不算，起點 $s$ 要算。 把兩個和併在一起：
\[
w(P) = \sum_{i=0}^{k-1} \bigl(ew[v_i, v_{i+1}] + nw[v_i]\bigr).
\]
每一條邊 $(v_i, v_{i+1})$ 正好帶著它\textbf{起點}的節點權重，而路徑上除了終點以外的每個節點，都恰好是
一條路徑邊的起點。 所以只要把每條邊的權重改成
\[
ew'[u, v] = ew[u, v] + nw[u],
\]
新定義下的最短路徑，就是邊權重為 $ew'$ 的圖上原本意義的最短路徑。 Dijkstra 本身不用改，只要鬆弛時
多加 $nw[u]$。

\textbf{擴充後的演算法}（改動的是 Input、第 6、7 行，另外補上第 9 行；$dist[v]$ 變小時要更新它在
$Q$ 裡的鍵，原題省略了這一步，這裡寫出來）：
\begin{center}
\small
\begin{tabular}{|p{0.92\linewidth}|}
\hline
\textbf{Algorithm}: Extended DSPA（考慮節點權重）\\ \hline
\textbf{Input}: $G=(V, E)$, $ew$, $nw$, $s$ \quad // 邊權重在 $ew$，節點權重在 $nw$，$s$ 是起點\\
\textbf{Output}: $SP$ \quad // 從 $s$ 到其他每個節點的最短路徑（終點的節點權重不計）\\ \hline
1:\ \ $dist[s]\leftarrow 0$; $dist[u]\leftarrow\infty$, for each $u\neq s$, $u\in V$\\
2:\ \ \textbf{insert} $u$ with key $dist[u]$ into priority queue $Q$, for each $u\in V$\\
3:\ \ \textbf{while} ($Q\neq\varnothing$)\\
4:\ \ \hspace*{1.5em}$u\leftarrow$Extract-Min($Q$)\\
5:\ \ \hspace*{1.5em}\textbf{for} each $v$ adjacent to $u$\\
6:\ \ \hspace*{3em}\textbf{if} $dist[v] > dist[u]+nw[u]+ew[u,v]$ \textbf{then}\\
7:\ \ \hspace*{4.5em}$dist[v]\leftarrow dist[u]+nw[u]+ew[u,v]$\\
8:\ \ \hspace*{4.5em}$pred[v]\leftarrow u$\quad // $v$ 在最短路徑上的前一個節點是 $u$\\
9:\ \ \hspace*{4.5em}Decrease-Key($Q$, $v$, $dist[v]$)\\
10: \textbf{calculate} the shortest path from $s$ to $u$ to add into set $SP$ according to $pred[u]$, for each $u\in V$, $u\neq s$\\
11: \textbf{return} $SP$\\ \hline
\end{tabular}
\end{center}

\textbf{正確性}。 上面說明了新的路徑權重等於邊權重 $ew'$ 下的路徑權重，第 6、7 行正是用 $ew'$ 做
鬆弛。 Dijkstra 的正確性只需要邊權重非負（資結書定理~25.6）：$ew[u, v] > 0$、$nw[u] \ge 0$，
所以 $ew'[u, v] > 0$，結論照樣成立，$dist[t]$ 就是新定義下 $s$ 到 $t$ 的最短距離，$pred$ 記下的
是對應的最短路徑。 $dist[s] = 0$ 也對：$s$ 到 $s$ 的路徑只有終點 $s$，不計權重。

如果節點權重可以是負的，只要每條邊都滿足 $ew[u, v] + nw[u] \ge 0$，上面的論證一樣成立；否則
$ew'$ 有負邊，Dijkstra 可能算錯（資結書第~25.2.5~節），要改用 Bellman--Ford。

\textbf{複雜度}。 每次鬆弛只多一次加法，和原來的 Dijkstra 一樣：用二元堆積是
$O((|V| + |E|) \lg |V|)$（資結書第~25.2~節）。

\textbf{小例子}。 $s \to a$ 權重 $1$、$a \to t$ 權重 $1$、$s \to t$ 權重 $5$，$nw[s] = 0$、
$nw[a] = 10$。 不算節點權重時最短路徑是 $s \to a \to t$，長 $2$；算進節點權重後它變成
$1 + 1 + 0 + 10 = 12$，最短路徑改成直接 $s \to t$，長 $5 + 0 = 5$。 擴充後的演算法：$s$ 取出時
$dist[a] = 1$、$dist[t] = 5$；$a$ 取出時檢查 $1 + 10 + 1 = 12 > 5$，不更新，$dist[t] = 5$，
$pred[t] = s$，正確。
\end{pastsoln}
